% Generated draft: requires mathematical and source review.
\documentclass[11pt]{article}
\usepackage[margin=25mm]{geometry}
\usepackage{fontspec}
\setmainfont{Latin Modern Roman}
\usepackage{amsmath,amssymb,mathtools,unicode-math}
\setmathfont{Latin Modern Math}
\usepackage{enumitem,needspace,xurl,hyperref}
\setlength{\parindent}{0pt}
\setlength{\parskip}{5pt}
\setlength{\emergencystretch}{3em}
\newcommand{\N}{\mathbb{N}}
\newcommand{\Z}{\mathbb{Z}}
\newcommand{\Q}{\mathbb{Q}}
\newcommand{\R}{\mathbb{R}}
\newcommand{\C}{\mathbb{C}}
\newcommand{\F}{\mathbb{F}}
\newcommand{\A}{\mathbb{A}}
\newcommand{\PP}{\mathbb{P}}
\newcommand{\E}{\mathbb{E}}
\newcommand{\eps}{\varepsilon}
\newcommand{\dd}{\,\mathrm{d}}
\begin{document}
\section*{UnsolvedMath ID: 30003692}
\textbf{Source problem code:} catalogue entry
\subsection{Definitions and mathematical statement}
The supplied record combines three independent questions and marks the combined statement as unrecoverable; it supplies no promoted clean statement. The questions are therefore retained separately, with unresolved notation and hypotheses made explicit.

\begin{enumerate}
\item For represented multivalued problems \(P,Q\), write \(P\leq_W Q\) for Weihrauch reducibility: computable preprocessing and postprocessing transform any realizer of \(Q\) into a realizer of \(P\), with the postprocessor also allowed access to the original input. Strictness means \(P\leq_W Q\) and \(Q\not\leq_W P\). The extracted question asks whether the displayed reduction
\[
\mathrm{List2N}\leq_W\mathrm{UCNN}
\]
is strict. The strings \(\mathrm{List2N}\) and \(\mathrm{UCNN}\) reproduce the supplied notation; their underlying problems, domains, representations, and intended subscripts or superscripts are not defined in the record. Consequently, this is only a conditional interpretation of strictness, not a complete specification of the two Weihrauch problems. The source also uses the plural ``reductions'' despite displaying only one inequality.

\item Is every computably compact, computably Hausdorff space computably regular? These are effective properties whose precise meanings require a specified representation or effective presentation and definitions of the relevant effective operations. The supplied record gives none of these conventions. The question is preserved as supplied, but its exact quantified class of spaces and effective separation requirements remain unresolved.

\item Does there exist a model \(M\models\mathrm{ZF}\) satisfying the following conditions internally? Let
\[
\mathcal D_T=\mathcal P(\N)/{\equiv_T},
\qquad
A\equiv_T B\ \Longleftrightarrow\ A\leq_T B\ \text{ and }\ B\leq_T A,
\]
where \(A\leq_T B\) means that membership in \(A\) is decidable by a Turing machine with oracle \(B\). Order the Turing degrees by \([A]_T\leq_T[B]_T\) if \(A\leq_T B\). There must exist \(C\subseteq\mathcal D_T\) and an order isomorphism
\[
f:(\omega_1,<)\longrightarrow(C,\leq_T)
\]
such that
\[
\forall d\in\mathcal D_T\ \exists c\in C\quad d\leq_T c
\]
and
\[
\forall d\in\mathcal D_T\quad
\left[\bigl(\forall c\in C\ (d\leq_T c\ \lor\ c\leq_T d)\bigr)
\Longrightarrow d\in C\right].
\]
The first condition says that \(C\) is cofinal; the second says that it is maximal among chains under inclusion. In addition, there must be no relation well-ordering the reals. Here reals are represented by subsets of \(\N\), and \(\omega_1\) is the least uncountable ordinal. All these sets, relations, quantifiers, and properties are interpreted inside \(M\), including \(\omega_1\) and the absence of a well-ordering. No countability, transitivity, or well-foundedness requirement on \(M\) is supplied.
\end{enumerate}

The supplied dated curation reports an affirmative answer to the first question and a counterexample to the assertion in the second. It further reports that the cited consistency result supplies a cofinal chain, rather than the cofinal maximal chain required in the third question, and labels that final question open. These are supplied literature assessments, not independently verified mathematical conclusions. In particular, cofinality must not replace maximality in the final question.
\subsection{Short English statement}
This record combines questions about a strict Weihrauch reduction, effective regularity of compact Hausdorff spaces, and a model of \(\mathrm{ZF}\) with a cofinal maximal chain of Turing degrees of order type \(\omega_1\) but no well-ordering of the reals. The first two questions lack essential effective definitions. Supplied dated curation reports that they were answered and that the third remains open; these reports require independent verification.
\subsection{Sources}
\begin{itemize}
\item Local UnsolvedMath registry: supplied complete record with permanent ID \(30003692\), source problem code catalogue entry, title ``Composite of Three Independent Computability Questions,'' in the Oberwolfach Reports Open Problems collection. No individual proposer is supplied. The proposed year is 2018.
\item Supplied source citation: \emph{Computability Theory (2019)}, Oberwolfach Reports, DOI 10.4171/owr/2018/1, pp.\ 5--41. \url{https://doi.org/10.4171/owr/2018/1}. The background describes Oberwolfach Reports as documenting mathematical workshops and their open-problem discussions.
\item Supplied literature source: Schroder, \emph{Effectively Hausdorff spaces}. \url{https://ems.press/content/serial-article-files/46899}.
\item Supplied literature source: Wang, Wu, and Yu, \emph{Cofinal maximal chains}. \url{https://doi.org/10.1090/S0002-9939-2014-11868-5}.
\item Literature assessment dated exactly 2026-08-21: the supplied curation labels the record open, reports that the first subquestion was answered yes and the compact-Hausdorff regularity assertion was disproved, and distinguishes a cited cofinal-chain consistency result from the requested cofinal maximal chain. It reports the final question as remaining open. This is supplied dated curation derived from the supplied OWR material, not an independently verified assessment; its time-sensitive status and cited claims require verification.
\item The supplied statement-verification note says that the official open-problem section contains independent questions, two answered during or before publication. The registry marks the extraction as unrecoverable, retains the original for provenance, and promotes no separate clean statement. These are supplied source-verification claims; no source visit is asserted here.
\end{itemize}
\end{document}
